% Chapter 2: Modules
\chapter{Modules}\label{chapter2}

One of the things which distinguishes the modern approach to Commutative Algebra is the greater emphasis on modules, rather than just on ideals. The extra ``elbow-room'' that this gives makes for greater clarity and simplicity. For instance, an ideal $\mathfrak{a}$ and its quotient ring $A/\mathfrak{a}$ are both examples of modules and so, to a certain extent, can be treated on an equal footing. In this chapter we give the definition and elementary properties of modules. We also give a brief treatment of tensor products, including a discussion of how they behave for exact sequences.

%=============================================================================
\section*{Modules and Module Homomorphisms}
\addcontentsline{toc}{section}{Modules and Module Homomorphisms}

Let $A$ be a ring (commutative, as always). An \emph{$A$-module}\index{module} is an abelian group $M$ (written additively) on which $A$ acts linearly: more precisely, it is a pair $(M,\mu)$, where $M$ is an abelian group and $\mu$ is a mapping of $A \times M$ into $M$ such that, if we write $ax$ for $\mu(a, x)$ ($a \in A$, $x \in M$), the following axioms are satisfied:
\begin{align*}
    a(x + y) &= ax + ay, \\
    (a + b)x &= ax + bx, \\
    (ab)x &= a(bx), \\
    1x &= x \qquad (a, b \in A;\ x, y \in M).
\end{align*}
(Equivalently, $M$ is an abelian group together with a ring homomorphism $A\to E(M)$, where $E(M)$ is the ring of endomorphisms of the abelian group $M$.)

\begin{examples}
1) An ideal $\mathfrak{a}$ of $A$ is an $A$-module. In particular $A$ itself is an $A$-module.

2) If $A$ is a field $k$, then $A$-module $=$ $k$-vector space.

3) $A = \Z$, then $\Z$-module $=$ abelian group (define $nx$ to be $x + \cdots + x$).

4) $A = k[x]$ where $k$ is a field; an $A$-module is a $k$-vector space with a linear transformation.

5) $G =$ finite group, $A = k[G] =$ group-algebra of $G$ over the field $k$ (thus $A$ is not commutative, unless $G$ is). Then $A$-module $=$ $k$-representation of $G$.
\end{examples}

Let $M, N$ be $A$-modules. A mapping $f \colon M \to N$ is an \emph{$A$-module homomorphism}\index{module homomorphism} (or is $A$-linear) if
\begin{align*}
    f(x + y) &= f(x) + f(y), \\
    f(ax) &= af(x)
\end{align*}
for all $a \in A$ and all $x, y \in M$. Thus $f$ is a homomorphism of abelian groups which commutes with the action of each $a\in A$. If $A$ is a field, an $A$-module homomorphism is the same thing as a linear transformation of vector spaces.

The composition of $A$-module homomorphisms is again an $A$-module homomorphism.

The set of all $A$-module homomorphisms from $M$ to $N$ is denoted by $\Hom_A(M, N)$ (or just $\Hom(M, N)$ if there is no ambiguity about the ring $A$). This is an $A$-module with addition and scalar multiplication defined by:
\begin{align*}
    (f + g)(x) &= f(x) + g(x), \\
    (af)(x) &= af(x)
\end{align*}
for all $x \in M$. It is a trivial matter to check that the axioms for an $A$-module are satisfied. This $A$-module is denoted by $\Hom_A(M,N)$ (or just $\Hom(M,N)$ if there is no ambiguity about the ring $A$).

Homomorphisms $u \colon M' \to M$ and $v \colon N \to N''$ induce mappings
\[
    \bar{u} \colon \Hom(M, N) \to \Hom(M', N) \qquad \text{and} \qquad \bar{v} \colon \Hom(M, N) \to \Hom(M, N'')
\]
defined as follows:
\[
    \bar{u}(f) = f \circ u, \qquad \bar{v}(f) = v \circ f.
\]
These mappings are $A$-module homomorphisms.

For any module $M$ there is a natural isomorphism $\Hom(A, M) \cong M$: any $A$-module homomorphism $f \colon A \to M$ is uniquely determined by $f(1)$, which can be any element of $M$.

%=============================================================================
\section*{Submodules and Quotient Modules}
\addcontentsline{toc}{section}{Submodules and Quotient Modules}

A \emph{submodule}\index{submodule} $M'$ of $M$ is a subgroup of $M$ which is closed under multiplication by elements of $A$. The abelian group $M/M'$ then inherits an $A$-module structure from $M$, defined by $a(x + M') = ax + M'$. The $A$-module $M/M'$ is the \emph{quotient}\index{quotient module} of $M$ by $M'$.

If $f \colon M \to N$ is an $A$-module homomorphism, the \emph{kernel}\index{kernel} of $f$ is
\[
    \Ker(f) = \{x \in M : f(x) = 0\}
\]
and is a submodule of $M$. The \emph{image}\index{image} of $f$ is
\[
    \Image(f) = f(M)
\]
and is a submodule of $N$. The \emph{cokernel}\index{cokernel} of $f$ is
\[
    \Coker(f) = N/\Image(f).
\]

If $M'$ is a submodule of $M$ such that $M' \subseteq \Ker(f)$, then $f$ gives rise to a homomorphism $\bar{f} \colon M/M' \to N$, defined as follows: if $\bar{x} \in M/M'$ is the image of $x \in M$, then $\bar{f}(\bar{x}) = f(x)$. The kernel of $\bar{f}$ is $\Ker(f)/M'$. In particular, taking $M' = \Ker(f)$, we have an isomorphism of $A$-modules
\[
    M/\Ker(f) \cong \Image(f).
\]

\section*{Operations on Submodules}
\addcontentsline{toc}{section}{Operations on Submodules}

Most of the operations on ideals considered in Chapter 1 have their counterparts for modules. Let $M$ be an $A$-module and let $(M_i)_{i \in I}$ be a family of submodules of $M$. Their \emph{sum} $\sum M_i$ is the set of all (finite) sums $\sum x_i$, where $x_i \in M_i$ for all $i \in I$, and almost all the $x_i$ (that is, all but a finite number) are zero. $\sum M_i$ is the smallest submodule of $M$ which contains all the $M_i$.

The intersection $\bigcap M_i$ is again a submodule of $M$. Thus the submodules of $M$ form a complete lattice with respect to inclusion.

\begin{proposition}\label{prop:submodule-isomorphisms}
\begin{enumerate}[i)]
    \item If $L \supseteq M \supseteq N$ are $A$-modules, then
    \[
        (L/N)/(M/N) \cong L/M.
    \]
    \item If $M_1, M_2$ are submodules of $M$, then
    \[
        (M_1 + M_2)/M_1 \cong M_2/(M_1 \cap M_2).
    \]
\end{enumerate}
\end{proposition}

\begin{proof}
i) Define $\theta \colon L/N \to L/M$ by $\theta(x + N) = x + M$. Then $\theta$ is a well-defined $A$-module homomorphism of $L/N$ onto $L/M$, and its kernel is $M/N$; hence (i).

ii) The composite homomorphism $M_2 \to M_1 + M_2 \to (M_1 + M_2)/M_1$ is surjective, and its kernel is $M_1 \cap M_2$; hence (ii).
\end{proof}

We cannot in general define the \emph{product} of two submodules, but we can define the product $\mathfrak{a}M$, where $\mathfrak{a}$ is an ideal and $M$ an $A$-module; it is the set of all finite sums $\sum a_i x_i$ with $a_i \in \mathfrak{a}$, $x_i \in M$, and is a submodule of $M$.

If $N, P$ are submodules of $M$, we define $(N : P)$ to be the set of all $a \in A$ such that $aP \subseteq N$; it is an \emph{ideal} of $A$. In particular, $(0 : M)$ is the set of all $a \in A$ such that $aM = 0$; this ideal is called the \emph{annihilator} of $M$ and is also denoted by $\operatorname{Ann}(M)$. If $\mathfrak{a} \subseteq \operatorname{Ann}(M)$, we may regard $M$ as an $A/\mathfrak{a}$-module, as follows: if $\bar{x} \in A/\mathfrak{a}$ is represented by $x \in A$, define $\bar{x}m$ to be $xm$ $(m \in M)$: this is independent of the choice of the representative $x$ of $\bar{x}$, since $\mathfrak{a}M = 0$.

An $A$-module is \emph{faithful}\index{faithful module} if $\Ann(M) = 0$. If $\Ann(M) = \mathfrak{a}$, then $M$ is faithful as an $A/\mathfrak{a}$-module.

\begin{exercise}\label{ex:annihilator-properties}
\begin{enumerate}[i)]
    \item $\Ann(M + N) = \Ann(M) \cap \Ann(N)$.
    \item $(N:P) = \Ann((N + P)/N)$.
\end{enumerate}
\end{exercise}

If $x$ is an element of $M$, the set of all multiples $ax$ $(a \in A)$ is a submodule of $M$, denoted by $Ax$ or $(x)$. If $M = \sum_{i \in I} Ax_i$, the $x_i$ are said to be a \emph{set of generators}\index{set of generators} of $M$; this means that every element of $M$ can be expressed (not necessarily uniquely) as a finite linear combination of the $x_i$ with coefficients in $A$. An $A$-module $M$ is said to be \emph{finitely generated}\index{finitely generated} if it has a finite set of generators.
%=============================================================================
\section*{Direct Sum and Product}
\addcontentsline{toc}{section}{Direct Sum and Product}

If $M, N$ are $A$-modules, their \emph{direct sum}\index{direct sum} $M \oplus N$ is the set of all pairs $(x, y)$ with $x \in M$, $y \in N$, with componentwise addition and scalar multiplication in the obvious way:
\begin{align*}
    (x_1, y_1) + (x_2, y_2) &= (x_1 + x_2, y_1 + y_2), \\
    a(x, y) &= (ax, ay).
\end{align*}

More generally, if $(M_i)_{i \in I}$ is any family of $A$-modules, we can define their \emph{direct sum}\index{direct sum} $\bigoplus_{i \in I} M_i$; its elements are families $(x_i)_{i \in I}$ such that $x_i \in M_i$ for each $i \in I$ and almost all $x_i$ are $0$. If we drop the restriction on the number of non-zero $x_i$'s we have the \emph{direct product}\index{direct product} $\prod_{i \in I} M_i$. Direct sum and direct product are therefore the same if the index set $I$ is finite, but not otherwise, in general.

Suppose that the ring $A$ is a direct product $\prod_{i=1}^n A_i$ (Chapter \ref{chapter1}). Then the set of all elements of $A$ of the form
\[
    (0, \ldots, 0, a_i, 0, \ldots, 0)
\]
with $a_i \in A_i$ is an \emph{ideal}\index{ideal} $\mathfrak{a}_i$ of $A$ (it is not a subring of A---except in trivial cases---because it does not contain the identity element of $A$). The ring $A$, considered as an $A$-module, is the direct sum of the ideals $\mathfrak{a}_1, \ldots, \mathfrak{a}_n$. Conversely, given a module decomposition
\[
    A = \mathfrak{a}_1 \oplus \cdots \oplus \mathfrak{a}_n
\]
of $A$ as a direct sum of ideals, we have
\[
    A \cong \prod_{i=1}^n A/\mathfrak{b}_i
\]
where $\mathfrak{b}_i = \bigoplus_{j \neq i} \mathfrak{a}_j$. Each ideal $\mathfrak{a}_i$ is a ring (isomorphic to $A/\mathfrak{b}_i$). The identity element $e_i$ of $\mathfrak{a}_i$ is an idempotent in $A$, and $\mathfrak{a}_i = (e_i)$.

%=============================================================================
\section*{Finitely Generated Modules}
\addcontentsline{toc}{section}{Finitely Generated Modules}

A \emph{free}\index{free module} $A$-module is one which is isomorphic to an $A$-module of the form $\bigoplus_{i \in I} M_i$, where each $M_i \cong A$ (as an $A$-module). A finitely generated free $A$-module is therefore isomorphic to $A \oplus \cdots \oplus A$ ($n$ summands), which is denoted by $A^n$.

\begin{proposition}\label{prop:finitely-generated-quotient}
$M$ is a finitely generated $A$-module $\iff$ $M$ is isomorphic to a quotient of $A^n$ for some integer $n \geq 0$.
\end{proposition}

\begin{proof}
    $\implies$: Let $x_1, \ldots, x_n$ generate $M$. Define $\varphi \colon A^n \to M$ by $\varphi(a_1, \ldots, a_n) = a_1 x_1 + \cdots + a_n x_n$. Then $\varphi$ is an $A$-module homomorphism onto $M$, and therefore $M \cong A^n/\Ker(\varphi)$.

    $\impliedby$: We have an $A$-module homomorphism $\varphi$ of $A^n$ onto $M$. If $e_i = (0, \ldots, 0, 1, 0, \ldots, 0)$ (the $1$ being in the $i$th place), then the $e_i$ ($1 \leq i \leq n$) generate $A^n$, hence the $\varphi(e_i)$ generate $M$.
\end{proof}

\begin{proposition}\label{prop:module-endomorphism-equation}
    Let $M$ be a finitely generated $A$-module, let $\mathfrak{a}$ be an ideal of $A$, and let $\varphi$ be an $A$-module endomorphism of $M$ such that $\varphi(M) \subseteq \mathfrak{a}M$. Then $\varphi$ satisfies an equation of the form
    \[
        \varphi^n + a_1 \varphi^{n-1} + \cdots + a_n = 0
    \]
    where the $a_i$ are in $\mathfrak{a}$.
\end{proposition}

\begin{proof}
    Let $x_1, \ldots, x_n$ be a set of generators of $M$. Then each $\varphi(x_i) \in \mathfrak{a}M$, so that we have say $\varphi(x_i) = \sum_{j=1}^n a_{ij}x_j$ ($1 \leq i \leq n$; $a_{ij} \in \mathfrak{a}$), i.e.,
    \[
        \sum_{j=1}^n (\delta_{ij}\varphi - a_{ij})x_j = 0
    \]
    where $\delta_{ij}$ is the Kronecker delta. By multiplying on the left by the adjoint of the matrix $(\delta_{ij}\varphi - a_{ij})$ it follows that $\det(\delta_{ij}\varphi - a_{ij})$ annihilates each $x_i$, hence is the zero endomorphism of $M$. Expanding out the determinant, we have an equation of the required form.
\end{proof}

\begin{corollary}\label{cor:cayley-hamilton}
    Let $M$ be a finitely generated $A$-module and let $\mathfrak{a}$ be an ideal of $A$ such that $\mathfrak{a}M = M$. Then there exists $x \equiv 1 \pmod{\mathfrak{a}}$ such that $xM = 0$.
\end{corollary}

\begin{proof}
    Take $\varphi = \text{identity}$, $x = 1 + a_1 + \cdots + a_n$ in \eqref{prop:module-endomorphism-equation}.
\end{proof}

\begin{proposition}[Nakayama's Lemma]\label{pro:nakayama}
    Let $M$ be a finitely generated $A$-module and $\mathfrak{a}$ an ideal of $A$ contained in the Jacobson radical $\jacrad$ of $A$. Then $\mathfrak{a}M = M$ implies $M = 0$.
\end{proposition}

\begin{proof}[First Proof]
    By \eqref{cor:cayley-hamilton} we have $xM = 0$ for some $x \equiv 1 \pmod{\mathfrak{a}}$. By Proposition~1.9, $x$ is a unit in $A$, hence $M = x^{-1}xM = 0$.
\end{proof}

\begin{proof}[Second Proof]
    Suppose $M \neq 0$, and let $u_1, \ldots, u_n$ be a minimal set of generators of $M$. Then $u_n \in \mathfrak{a}M$, hence we have an equation of the form $u_n = a_1 u_1 + \cdots + a_n u_n$, with the $a_i \in \mathfrak{a}$. Hence
    \[
        (1 - a_n)u_n = a_1 u_1 + \cdots + a_{n-1}u_{n-1};
    \]
    since $a_n \in \jacrad$, it follows from \eqref{prop:jacobson-radical} that $1 - a_n$ is a unit in $A$. Hence $u_n$ belongs to the submodule of $M$ generated by $u_1, \ldots, u_{n-1}$: contradiction.
\end{proof}

\begin{corollary}\label{cor:nakayama-submodule}
    Let $M$ be a finitely generated $A$-module, $N$ a submodule of $M$, $\mathfrak{a} \subseteq \jacrad$ an ideal. Then $M = \mathfrak{a}M + N \implies M = N$.
\end{corollary}

\begin{proof}
    Apply \eqref{pro:nakayama} to $M/N$, observing that $\mathfrak{a}(M/N) = (\mathfrak{a}M + N)/N$.
\end{proof}

Let $A$ be a local ring, $\mathfrak{m}$ its maximal ideal, $k = A/\mathfrak{m}$ its residue field. Let $M$ be a finitely generated $A$-module. $M/\mathfrak{m}M$ is annihilated by $\mathfrak{m}$, hence is naturally an $A/\mathfrak{m}$-module, i.e., a $k$-vector space, and as such is finite-dimensional.

\begin{proposition}\label{prop:generate-modulo-maximal}
Let $x_i$ ($1 \leq i \leq n$) be elements of $M$ whose images in $M/\mathfrak{m}M$ form a basis of this vector space. Then the $x_i$ generate $M$.
\end{proposition}

\begin{proof}
Let $N$ be the submodule of $M$ generated by the $x_i$. Then the composite map $N \to M \to M/\mathfrak{m}M$ maps $N$ onto $M/\mathfrak{m}M$, hence $N + \mathfrak{m}M = M$, hence $N = M$ by \eqref{cor:nakayama-submodule}.
\end{proof}

%=============================================================================
\section*{Exact Sequences}
\addcontentsline{toc}{section}{Exact Sequences}

A sequence of $A$-modules and $A$-homomorphisms
\[
    \cdots \to M_{i-1} \stackrel{f_i}{\longrightarrow} M_i \stackrel{f_{i+1}}{\longrightarrow} M_{i+1} \longrightarrow \cdots \tag{0}
\]
is said to be \emph{exact}\index{exact sequence} at $M_i$ if $\Image(f_i) = \Ker(f_{i+1})$. The sequence is exact if it is exact at each $M_i$. In particular:

$0 \longrightarrow M' \stackrel{f}{\longrightarrow} M$ is exact $\iff f$ is injective; \hfill (1)

$M \stackrel{g}{\longrightarrow} M'' \longrightarrow 0$ is exact $\iff g$ is surjective; \hfill (2)

$0 \longrightarrow M' \stackrel{f}{\longrightarrow} M \stackrel{g}{\longrightarrow} M'' \longrightarrow 0$ is exact $\iff f$ is injective, $g$ is surjective and $g$ induces an isomorphism of $\operatorname{Coker}(f) = M/f(M')$ onto $M''$. \hfill (3)

A sequence of type (3) is called a \emph{short exact sequence}\index{short exact sequence}. Any long exact sequence (0) can be split up into short exact sequences: if $N_i = \Image(f_i) = \Ker(f_{i+1})$, we have short exact sequences $0 \longrightarrow N_i \longrightarrow M_i \longrightarrow N_{i+1} \longrightarrow 0$ for each $i$.

\begin{proposition}\label{prop:exact-hom}
\begin{enumerate}[i)]
    \item Let \[M' \stackrel{u}{\longrightarrow} M \stackrel{v}{\longrightarrow} M'' \longrightarrow 0 \tag{4}\] be a sequence. Then it is exact $\iff$ for all $A$-modules $N$, the sequence
    \[
        0 \longrightarrow \Hom(M'', N) \longrightarrow \Hom(M, N) \longrightarrow \Hom(M', N) \tag{4'}
    \]
    is exact.
    \item Let \[0 \longrightarrow N' \stackrel{u}{\longrightarrow} N \stackrel{v}{\longrightarrow} N''\tag{5}\] be a sequence. Then it is exact $\iff$ for all $A$-modules $M$, the sequence
    \[
        0 \longrightarrow \Hom(M, N') \longrightarrow \Hom(M, N) \longrightarrow \Hom(M, N'')\tag{5'}
    \]
    is exact.
\end{enumerate}
\end{proposition}

\begin{proof}
All four parts of this proposition are easy exercises. For example, suppose that (i) is exact for all $N$. First of all, since $\bar{v}$ is injective for all $N$ it follows that $v$ is surjective. Next, we have $\bar{u} \circ \bar{v} = 0$, that is $\bar{v} \circ \bar{u} = v \circ u \circ f = 0$ for all $f \colon M'' \to N$. Taking $N$ to be $M''$ and $f$ to be the identity mapping, it follows that $v \circ u = 0$, hence $\Image(u) \subseteq \Ker(v)$. Next take $N = M/\Image(u)$ and let $\phi \colon M \to N$ be the projection. Then $\phi \in \Ker(\bar{u})$, hence there exists $\psi \colon M'' \to N$ such that $\phi = \psi \circ v$. Consequently $\Image(u) = \Ker(\phi) \supseteq \Ker(v)$.
\end{proof}

\begin{proposition}[Snake Lemma]\label{prop:snake-lemma}
Let
\[\begin{tikzcd}
	0 & {M'} & M & {M''} & 0 \\
	0 & {N'} & N & {N''} & 0
	\arrow[from=1-1, to=1-2]
	\arrow["u", from=1-2, to=1-3]
	\arrow["{f'}"', from=1-2, to=2-2]
	\arrow["v", from=1-3, to=1-4]
	\arrow["f"', from=1-3, to=2-3]
	\arrow[from=1-4, to=1-5]
	\arrow["{f''}"', from=1-4, to=2-4]
	\arrow[from=2-1, to=2-2]
	\arrow["{u'}"', from=2-2, to=2-3]
	\arrow["{v'}"', from=2-3, to=2-4]
	\arrow[from=2-4, to=2-5]
\end{tikzcd}\]
be a commutative diagram of $A$-modules and homomorphisms, with the rows exact. Then there exists an exact sequence
\[
    0 \longrightarrow \Ker(f') \longrightarrow \Ker(f) \longrightarrow \Ker(f'') \stackrel{d}{\longrightarrow} \Coker(f') \longrightarrow \Coker(f) \longrightarrow \Coker(f'') \longrightarrow 0
\]
in which the first two maps are restrictions of $u, v$, and the last two are induced by $u', v'$.
\end{proposition}

\begin{proof}
    The boundary homomorphism $d$ is defined as follows: if $x'' \in \Ker(f'')$, we have $x'' = v(x)$ for some $x \in M$, and $v'(f(x)) = f''(v(x)) = 0$, hence $f(x) \in \Ker(v') = \Image(u')$, so that $f(x) = u'(y')$ for some $y' \in N'$. Then $d(x'')$ is defined to be the image of $y'$ in $\Coker(f')$. The verification that $d$ is well-defined, and that the sequence is exact, is a straightforward exercise in diagram-chasing which we leave to the reader.
\end{proof}

\begin{remark}
    \eqref{prop:snake-lemma} is a special case of the exact homology sequence of homological algebra.
\end{remark}

Let $C$ be a class of $A$-modules and let $\lambda$ be a function on $C$ with values in $\Z$ (or, more generally, with values in an abelian group $G$). The function $\lambda$ is \emph{additive}\index{additive function} if, for each short exact sequence
\[
    0 \longrightarrow M' \longrightarrow M \longrightarrow M'' \longrightarrow 0
\]
in which all the terms belong to $\mathcal{C}$, we have
\[
    \lambda(M') - \lambda(M) + \lambda(M'') = 0.
\]

\begin{example}
Let $A$ be a field $k$, and let $\mathcal{C}$ be the class of all finite-dimensional $k$-vector spaces $V$. Then $V \mapsto \dim V$ is an additive function on $\mathcal{C}$.
\end{example}

\begin{proposition}\label{prop:additive-function}
Let $0 \longrightarrow M_0 \longrightarrow M_1 \longrightarrow \cdots \longrightarrow M_n \longrightarrow 0$ be an exact sequence of $A$-modules in which all the modules $M_i$ and the kernels of all the homomorphisms belong to $\mathcal{C}$. Then for any additive function $\lambda$ on $\mathcal{C}$ we have
\[
    \sum_{i=0}^n (-1)^i \lambda(M_i) = 0.
\]
\end{proposition}

\begin{proof}
Split up the sequence into short exact sequences
\[
    0 \longrightarrow N_j \longrightarrow M_j \longrightarrow N_{j+1} \longrightarrow 0
\]
($N_0 = N_{n+1} = 0$). Then we have $\lambda(M_j) = \lambda(N_j) + \lambda(N_{j+1})$. Now take the alternating sum of the $\lambda(M_j)$, and everything cancels out.
\end{proof}

%=============================================================================
\section*{Tensor Product of Modules}
\addcontentsline{toc}{section}{Tensor Product of Modules}

Let $M, N, P$ be three $A$-modules. A mapping $f \colon M \times N \to P$ is said to be $A$-bilinear if for each $x \in M$ the mapping $y \mapsto f(x, y)$ of $N$ into $P$ is $A$-linear, and for each $y \in N$ the mapping $x \mapsto f(x, y)$ of $M$ into $P$ is $A$-linear.

We shall construct an $A$-module $T$, called the \textit{tensor product} of $M$ and $N$, with the property that the $A$-bilinear mappings $M \times N \to P$ are in a natural one-to-one correspondence with the $A$-linear mappings $T \to P$, for all $A$-modules $P$. More precisely:

\begin{proposition}\label{prop:tensor-universal}
Let $M, N$ be $A$-modules. Then there exists a pair $(T, g)$ consisting of an $A$-module $T$ and an $A$-bilinear mapping $g \colon M \times N \to T$, with the following property:

Given any $A$-module $P$ and any $A$-bilinear mapping $f \colon M \times N \to P$, there exists a unique $A$-linear mapping $f' \colon T \to P$ such that $f = f' \circ g$.

Moreover, if $(T, g)$ and $(T', g')$ are two pairs with this property, then there exists a unique isomorphism $j \colon T \to T'$ such that $j \circ g = g'$.
\end{proposition}

\begin{proof}
i) \textit{Uniqueness.} Replacing $(P, f)$ by $(T', g')$ we get a unique $j \colon T \to T'$ such that $g' = j \circ g$. Interchanging the roles of $T$ and $T'$, we get $j' \colon T' \to T$ such that $g = j' \circ g'$. Each of the compositions $j \circ j'$, $j' \circ j$ must be the identity, and therefore $j$ is an isomorphism.

ii) \textit{Existence.} Let $C$ denote the free $A$-module $A^{(M \times N)}$. The elements of $C$ are formal linear combinations of elements of $M \times N$ with coefficients in $A$, i.e.\ they are expressions of the form $\sum_{i=1}^r a_i \cdot (x_i, y_i)$ ($a_i \in A$, $x_i \in M$, $y_i \in N$).

Let $D$ be the submodule of $C$ generated by all elements of $C$ of the following types:
\begin{align*}
    &(x + x', y) - (x, y) - (x', y), \\
    &(x, y + y') - (x, y) - (x, y'), \\
    &(ax, y) - a \cdot (x, y), \\
    &(x, ay) - a \cdot (x, y).
\end{align*}

Let $T = C/D$. For each basis element $(x, y)$ of $C$, let $x \otimes y$ denote its image in $T$. Then $T$ is generated by the elements of the form $x \otimes y$, and from our definitions we have
\begin{align*}
    (x + x') \otimes y &= x \otimes y + x' \otimes y, \\
    x \otimes (y + y') &= x \otimes y + x \otimes y', \\
    (ax) \otimes y &= x \otimes (ay) = a(x \otimes y).
\end{align*}
Equivalently, the mapping $g \colon M \times N \to T$ defined by $g(x, y) = x \otimes y$ is $A$-bilinear.

Any map $f$ of $M \times N$ into an $A$-module $P$ extends by linearity to an $A$-module homomorphism $\tilde{f} \colon C \to P$. Suppose in particular that $f$ is $A$-bilinear. Then, from the definitions, $\tilde{f}$ vanishes on all the generators of $D$, hence on the whole of $D$, and therefore induces a well-defined $A$-homomorphism $f'$ of $T = C/D$ into $P$ such that $f'(x \otimes y) = f(x, y)$. The mapping $f'$ is uniquely defined by this condition, and therefore the pair $(T, g)$ satisfy the conditions of the proposition.
\end{proof}

\begin{remark}\leavevmode
\begin{enumerate}[i)]
    \item The module $T$ constructed above is called the \emph{tensor product} of $M$ and $N$, and is denoted by $M \otimes_A N$, or just $M \otimes N$ if there is no ambiguity about the ring $A$. It is generated as an $A$-module by the ``products'' $x \otimes y$. If $(x_i)_{i \in I}$, $(y_j)_{j \in J}$ are families of generators of $M, N$ respectively, then the elements $x_i \otimes y_j$ generate $M \otimes N$. In particular, if $M$ and $N$ are finitely generated, so is $M \otimes N$.
    
    \item The notation $x \otimes y$ is inherently ambiguous unless we specify the tensor product to which it belongs. Let $M', N'$ be submodules of $M, N$ respectively, and let $x \in M'$ and $y \in N'$. Then it can happen that $x \otimes y$ as an element of $M \otimes N$ is zero whilst $x \otimes y$ as an element of $M' \otimes N'$ is non-zero. For example, take $A = \Z$, $M = \Z$, $N = \Z/2\Z$, and let $M'$ be the submodule $2\Z$ of $\Z$, whilst $N' = N$. Let $x$ be the non-zero element of $N$ and consider $2 \otimes x$. As an element of $M \otimes N$, it is zero because $2 \otimes x = 1 \otimes 2x = 1 \otimes 0 = 0$. But as an element of $M' \otimes N'$ it is non-zero. See the example after \eqref{prop:tensor-right-exact}.
    
    However, there is the following result:

    \begin{corollary}
        Let $x_i \in M$, $y_i \in N$ be such that $\sum x_i \otimes y_i = 0$ in $M \otimes N$. Then there exist finitely generated submodules $M_0$ of $M$ and $N_0$ of $N$ such that $\sum x_i \otimes y_i = 0$ in $M_0 \otimes N_0$.
    \end{corollary}

    \begin{proof}
        If $\sum x_i \otimes y_i = 0$ in $M \otimes N$, then in the notation of the proof of \eqref{prop:tensor-universal} we have $\sum (x_i, y_i) \in D$, and therefore $\sum (x_i, y_i)$ is a finite sum of generators of $D$. Let $M_0$ be the submodule of $M$ generated by the $x_i$ and all the elements of $M$ which occur as first coordinates in these generators of $D$, and define $N_0$ similarly. Then $\sum x_i \otimes y_i = 0$ as an element of $M_0 \otimes N_0$.
    \end{proof}

    \item We shall never again need to use the construction of the tensor product given in \eqref{prop:tensor-universal}, and the reader may safely forget it if he prefers. What is essential to keep in mind is the defining property of the tensor product.
    \item Instead of starting with bilinear mappings we could have started with multilinear mappings $f: M_1 \times \cdots \times M_r \to P$ defined in the same way (i.e., linear in each variable). Following through the proof of \eqref{prop:tensor-universal} we should end up with a ``multi-tensor product'' $T = M_1 \otimes \cdots \otimes M_r$, generated by all products $x_1 \otimes \cdots \otimes x_r$ ($x_i \in M_i$, $1 \leq i \leq r$). The details may safely be left to the reader; the result corresponding to \eqref{prop:tensor-universal} is
\end{enumerate}
\end{remark}
\bgroup
    % 保存原始定义
    \let\savedtheproposition\theproposition
    \let\savedtheHproposition\theHproposition   % 保存锚点标识宏
    
    % 计数器退回到 12
    \addtocounter{proposition}{-2}      % 假设原本是 14
    
    % 修改显示格式：2.12*
    \renewcommand{\theproposition}{\savedtheproposition*}
    % 修改锚点标识：使锚点唯一，比如加上 star 后缀
    \renewcommand{\theHproposition}{\savedtheHproposition.star} 
  \begin{proposition}\label{prop:multilinear-tensor}
    Let $M_1, \dots, M_r$ be $A$-modules. Then there exists a pair $(T, g)$ consisting of an $A$-module $T$ and an $A$-multilinear mapping $g: M_1 \times \cdots \times M_r \to T$ with the following property:

        Given any $A$-module $P$ and any $A$-multilinear mapping $f: M_1 \times \cdots \times M_r \to P$, there exists a unique $A$-homomorphism $f': T \to P$ such that $f' \circ g = f$.

        Moreover, if $(T, g)$ and $(T', g')$ are two pairs with this property, then there exists a unique isomorphism $j: T \to T'$ such that $j \circ g = g'$.
    \end{proposition}
\egroup
\addtocounter{proposition}{1}% 环境内计数器+1 变成13，再 +1 回到14


There are various so-called ``canonical isomorphisms'', some of which we state here:

\begin{proposition}\label{pro:canonical-isomorphisms}
Let $M, N, P$ be $A$-modules. Then there exist unique isomorphisms
\begin{enumerate}[i)]
    \item $M \otimes N \to N \otimes M$
    \item $(M \otimes N) \otimes P \to M \otimes (N \otimes P) \to M \otimes N \otimes P$
    \item $(M \oplus N) \otimes P \to (M \otimes P) \oplus (N \otimes P)$
    \item $A \otimes M \to M$
\end{enumerate}
such that, respectively,
\begin{enumerate}
    \item[a)] $x \otimes y \mapsto y \otimes x$
    \item[b)] $(x \otimes y) \otimes z \mapsto x \otimes (y \otimes z) \mapsto x \otimes y \otimes z$
    \item[c)] $(x, y) \otimes z \mapsto (x \otimes z, y \otimes z)$
    \item[d)] $a \otimes x \mapsto ax$.
\end{enumerate}
\end{proposition}

\begin{proof}
In each case the point is to show that the mappings so described are well defined. The technique is to construct suitable bilinear or multilinear mappings, and use the defining property \eqref{prop:tensor-universal} or \eqref{prop:multilinear-tensor} to infer the existence of homomorphisms of tensor products. We shall prove half of ii) as an example of the method, and leave the rest to the reader.

We shall construct homomorphisms
\[
(M \otimes N) \otimes P \stackrel{f}{\longrightarrow} M \otimes N \otimes P \stackrel{g}{\longrightarrow} (M \otimes N) \otimes P
\]
such that $f((x \otimes y) \otimes z) = x \otimes y \otimes z$ and $g(x \otimes y \otimes z) = (x \otimes y) \otimes z$ for all $x \in M$, $y \in N$, $z \in P$.

To construct $f$, fix the element $z \in P$. The mapping $(x, y) \mapsto x \otimes y \otimes z$ ($x \in M$, $y \in N$) is bilinear in $x$ and $y$ and therefore induces a homomorphism $f_z: M \otimes N \to M \otimes N \otimes P$ such that $f_z(x \otimes y) = x \otimes y \otimes z$. Next, consider the mapping $(t, z) \mapsto f_z(t)$ of $(M \otimes N) \times P$ into $M \otimes N \otimes P$. This is bilinear in $t$ and $z$ and therefore induces a homomorphism
\[
f: (M \otimes N) \otimes P \to M \otimes N \otimes P
\]
such that $f((x \otimes y) \otimes z) = x \otimes y \otimes z$.

To construct $g$, consider the mapping $(x, y, z) \mapsto (x \otimes y) \otimes z$ of $M \times N \times P$ into $(M \otimes N) \otimes P$. This is linear in each variable and therefore induces a homomorphism
\[
g: M \otimes N \otimes P \to (M \otimes N) \otimes P
\]
such that $g(x \otimes y \otimes z) = (x \otimes y) \otimes z$.

Clearly $f \circ g$ and $g \circ f$ are identity maps, hence $f$ and $g$ are isomorphisms. 
\end{proof}

\begin{exercise}\label{ex:bimodule-tensor}
Let $A, B$ be rings, let $M$ be an $A$-module, $P$ a $B$-module and $N$ an $(A, B)$-bimodule (that is, $N$ is simultaneously an $A$-module and a $B$-module and the two structures are compatible in the sense that $a(xb) = (ax)b$ for all $a \in A$, $b \in B$, $x \in N$). Then $M \otimes_A N$ is naturally a $B$-module, $N \otimes_B P$ an $A$-module, and we have
\[
    (M \otimes_A N) \otimes_B P \cong M \otimes_A (N \otimes_B P).
\]
\end{exercise}

Let $f \colon M \to M'$, $g \colon N \to N'$ be homomorphisms of $A$-modules. Define $h \colon M \times N \to M' \otimes N'$ by $h(x, y) = f(x) \otimes g(y)$. It is easily checked that $h$ is $A$-bilinear and therefore induces an $A$-module homomorphism
\[
    f \otimes g \colon M \otimes N \to M' \otimes N'
\]
such that
\[
    (f \otimes g)(x \otimes y) = f(x) \otimes g(y) \qquad (x \in M, \quad y \in N).
\]

Let $f' \colon M' \to M''$ and $g' \colon N' \to N''$ be homomorphisms of $A$-modules. Then clearly the homomorphisms $(f' \circ f) \otimes (g' \circ g)$ and $(f' \otimes g') \circ (f \otimes g)$ agree on all elements of the form $x \otimes y$ in $M \otimes N$. Since these elements generate $M \otimes N$, it follows that
\[
    (f' \circ f) \otimes (g' \circ g) = (f' \otimes g') \circ (f \otimes g).
\]

\section*{Restriction and extension of scalars}
\addcontentsline{toc}{section}{Restriction and extension of scalars}

Let $f \colon A \to B$ be a homomorphism of rings and let $N$ be a $B$-module. Then $N$ has an $A$-module structure defined as follows: if $a \in A$ and $x \in N$, then $ax$ is defined to be $f(a)x$. This $A$-module is said to be obtained from $N$ by \emph{restriction of scalars}\index{restriction of scalars}. In particular, $f$ defines in this way an $A$-module structure on $B$.

\begin{proposition}\label{prop:finitely-generated-restriction}
Suppose $N$ is finitely generated as a $B$-module and that $B$ is finitely generated as an $A$-module. Then $N$ is finitely generated as an $A$-module.
\end{proposition}

\begin{proof}
Let $y_1, \ldots, y_n$ generate $N$ over $B$, and let $x_1, \ldots, x_m$ generate $B$ as an $A$-module. Then the $mn$ products $x_i y_j$ generate $N$ over $A$.
\end{proof}

Now let $M$ be an $A$-module. Since, as we have just seen, $B$ can be regarded as an $A$-module, we can form the $A$-module $M_B = B \otimes_A M$. In fact $M_B$ carries a $B$-module structure such that $b(b' \otimes x) = bb' \otimes x$ for all $b, b' \in B$ and all $x \in M$. The $B$-module $M_B$ is said to be obtained from $M$ by \emph{extension of scalars}\index{extension of scalars}.

\begin{proposition}\label{prop:finitely-generated-extension}
If $M$ is finitely generated as an $A$-module, then $M_B$ is finitely generated as a $B$-module.
\end{proposition}

\begin{proof}
If $x_1, \ldots, x_m$ generate $M$ over $A$, then the $1 \otimes x_i$ generate $M_B$ over $B$.
\end{proof}

%=============================================================================
\section*{Exactness Properties of the Tensor Product}
\addcontentsline{toc}{section}{Exactness Properties of the Tensor Product}

Let $f \colon M \times N \to P$ be an $A$-bilinear mapping. For each $x \in M$ the mapping $y \mapsto f(x, y)$ of $N$ into $P$ is $A$-linear, hence $f$ gives rise to a mapping $M \to \Hom(N, P)$ which is $A$-linear because $f$ is linear in the variable $x$. Conversely any $A$-homomorphism $\phi \colon M \to \Hom_A(N, P)$ defines a bilinear map, namely $(x, y) \mapsto \phi(x)(y)$. Hence the set $S$ of $A$-bilinear mappings $M \times N \to P$ is in natural one-to-one correspondence with $\Hom(M, \Hom(N, P))$. On the other hand $S$ is in one-to-one correspondence with $\Hom(M \otimes N, P)$, by the defining property of the tensor product. Hence we have a canonical isomorphism
\begin{equation}\label{eq:tensor-hom-adjunction}
    \Hom(M \otimes N, P) \cong \Hom(M, \Hom(N, P)). \tag{1}
\end{equation}

\begin{proposition}\label{prop:tensor-right-exact}
Let \begin{equation}\label{eq:exact-sequence}
    M' \stackrel{f}{\longrightarrow} M \stackrel{g}{\longrightarrow} M'' \longrightarrow 0 \tag{2}
\end{equation}
be an exact sequence of $A$-modules and homomorphisms, and let $N$ be any $A$-module. Then the sequence
\begin{equation}\label{eq:tensored-sequence}
    M' \otimes N \xrightarrow{f \otimes 1} M \otimes N \xrightarrow{g \otimes 1} M'' \otimes N \longrightarrow 0 \tag{3}
\end{equation}
(where $1$ denotes the identity mapping on $N$) is exact.
\end{proposition}

\begin{proof}
Let $E$ denote the sequence \eqref{eq:exact-sequence}, and let $E \otimes N$ denote the tensored sequence \eqref{eq:tensored-sequence}. Let $P$ be any $A$-module. Since \eqref{eq:exact-sequence} is exact, the sequence $\Hom(E, \Hom(N, P))$ is exact by \eqref{prop:exact-hom}; hence by \eqref{eq:tensor-hom-adjunction} the sequence $\Hom(E \otimes N, P)$ is exact. By \eqref{prop:exact-hom} again, it follows that $E \otimes N$ is exact.
\end{proof}

\begin{remarks}
    \begin{enumerate}[i)]
        \item Let $T(M) = M \otimes N$ and let $U(P) = \Hom(N, P)$. Then \eqref{eq:tensor-hom-adjunction} takes the form $\Hom(T(M), P) = \Hom(M, U(P))$ for all $A$-modules $M$ and $P$. In the language of abstract nonsense, the functor $T$ is the left adjoint of $U$, and $U$ is the right adjoint of $T$. The proof of \eqref{prop:tensor-right-exact} shows that any functor which is a left adjoint is right exact. Likewise any functor which is a right adjoint is left exact.
        \item It is \emph{not} in general true that, if $M' \longrightarrow M \longrightarrow M''$ is an exact sequence of $A$-modules and homomorphisms, the sequence $M' \otimes N \longrightarrow M \otimes N \longrightarrow M'' \otimes N$ obtained by tensoring with an arbitrary $A$-module $N$ is exact.
    \end{enumerate}
\end{remarks}

\begin{example}
Take $A = \Z$ and consider the exact sequence $0 \longrightarrow \Z \stackrel{f}{\longrightarrow} \Z$, where $f(x) = 2x$ for all $x \in \Z$. If we tensor with $N = \Z/2\Z$, the sequence $0 \longrightarrow \Z \otimes N \xrightarrow{f \otimes 1} \Z \otimes N$ is not exact, because for any $x \otimes y \in \Z \otimes N$ we have
\[
    (f \otimes 1)(x \otimes y) = 2x \otimes y = x \otimes 2y = x \otimes 0 = 0,
\]
so that $f \otimes 1$ is the zero mapping, whereas $\Z \otimes N \neq 0$.
\end{example}

The functor $T_N \colon M \mapsto M \otimes_A N$ on the category of $A$-modules and homomorphisms is therefore not in general exact. If $T_N$ is exact, that is to say if tensoring with $N$ transforms all exact sequences into exact sequences, then $N$ is said to be a \emph{flat}\index{flat module} $A$-module.

\begin{proposition}\label{prop:flat-equivalent}
The following are equivalent, for an $A$-module $N$:
\begin{enumerate}[i)]
    \item $N$ is flat.
    \item If $0 \longrightarrow M' \longrightarrow M \longrightarrow M'' \longrightarrow 0$ is any exact sequence of $A$-modules, the tensored sequence $0 \longrightarrow M' \otimes N \longrightarrow M \otimes N \longrightarrow M'' \otimes N \longrightarrow 0$ is exact.
    \item If $f \colon M' \to M$ is injective, then $f \otimes 1 \colon M' \otimes N \longrightarrow M \otimes N$ is injective.
    \item If $f \colon M' \to M$ is injective and $M, M'$ are finitely generated, then $f \otimes 1 \colon M' \otimes N \to M \otimes N$ is injective.
\end{enumerate}
\end{proposition}

\begin{proof}
i) $\iff$ ii) by splitting up a long exact sequence into short exact sequences.

ii) $\iff$ iii) by \eqref{prop:tensor-right-exact}.

iii) $\implies$ iv): clear.

iv) $\implies$ iii). Let $f \colon M' \to M$ be injective and let $u = \sum x_i \otimes y_i \in \Ker(f \otimes 1)$, so that $\sum f(x_i) \otimes y_i = 0$ in $M \otimes N$. Let $M'_0$ be the submodule of $M'$ generated by the $x_i$ and let $u_0$ denote $\sum x_i \otimes y_i$ as an element of $M'_0 \otimes N$. By \eqref{pro:canonical-isomorphisms} there exists a finitely generated submodule $M_0$ of $M$ containing $f(M'_0)$ and such that $\sum f(x_i) \otimes y_i = 0$ as an element of $M_0 \otimes N$. If $f_0 \colon M'_0 \to M_0$ is the restriction of $f$, this means that $(f_0 \otimes 1)(u_0) = 0$. Since $M_0$ and $M'_0$ are finitely generated, $f_0 \otimes 1$ is injective and therefore $u_0 = 0$, hence $u = 0$.
\end{proof}

\begin{exercise}\label{ex:flat-base-change}
If $f \colon A \to B$ is a ring homomorphism and $M$ is a flat $A$-module, then $M_B = B \otimes_A M$ is a flat $B$-module. (Use the canonical isomorphisms \eqref{pro:canonical-isomorphisms} and \eqref{ex:bimodule-tensor}.)
\end{exercise}

%=============================================================================
\section*{Algebras}
\addcontentsline{toc}{section}{Algebras}

Let $f \colon A \to B$ be a ring homomorphism. If $a \in A$ and $b \in B$, define a product
\[
    ab = f(a)b.
\]
This definition of scalar multiplication makes the ring $B$ into an $A$-module (it is a particular example of restriction of scalars). Thus $B$ has an $A$-module structure as well as a ring structure, and these two structures are compatible in a sense which the reader will be able to formulate for himself. The ring $B$, equipped with this $A$-module structure, is said to be an \emph{$A$-algebra}\index{algebra}. Thus an $A$-algebra is, by definition, a ring $B$ together with a ring homomorphism $f \colon A \to B$.

\begin{remark}\leavevmode
\begin{enumerate}[i)]
    \item In particular, if $A$ is a field $K$ (and $B \neq 0$) then $f$ is injective by Proposition~1.2 and therefore $K$ can be canonically identified with its image in $B$. Thus a $K$-algebra ($K$ a field) is effectively a ring containing $K$ as a subring.
    
    \item Let $A$ be any ring. Since $A$ has an identity element there is a unique homomorphism of the ring of integers $\Z$ into $A$, namely $n \mapsto n \cdot 1$. Thus every ring is automatically a $\Z$-algebra.
\end{enumerate}
\end{remark}

Let $f \colon A \to B$, $g \colon A \to C$ be two ring homomorphisms. An \emph{$A$-algebra homomorphism} $h \colon B \to C$ is a ring homomorphism which is also an $A$-module homomorphism. The reader should verify that $h$ is an $A$-algebra homomorphism if and only if $h \circ f = g$.

A ring homomorphism $f \colon A \to B$ is \emph{finite}\index{finite homomorphism}, and $B$ is a \emph{finite $A$-algebra}, if $B$ is finitely generated as an $A$-module. The homomorphism $f$ is of \emph{finite type}\index{finite type}, and $B$ is a \emph{finitely-generated $A$-algebra}, if there exists a finite set of elements $x_1, \ldots, x_n$ in $B$ such that every element of $B$ can be written as a polynomial in $x_1, \ldots, x_n$ with coefficients in $f(A)$; or equivalently if there is an $A$-algebra homomorphism from a polynomial ring $A[t_1, \ldots, t_n]$ onto $B$.

A ring $A$ is said to be \emph{finitely generated}\index{finitely generated ring} if it is finitely generated as a $\Z$-algebra. This means that there exist finitely many elements $x_1, \ldots, x_n$ in $A$ such that every element of $A$ can be written as a polynomial in the $x_i$ with rational integer coefficients.

%=============================================================================
\section*{Tensor Product of Algebras}
\addcontentsline{toc}{section}{Tensor Product of Algebras}

Let $B, C$ be two $A$-algebras, $f \colon A \to B$, $g \colon A \to C$ the corresponding homomorphisms. Since $B$ and $C$ are $A$-modules we may form their tensor product $D = B \otimes_A C$, which is an $A$-module. We shall now define a multiplication on $D$.

Consider the mapping $B \times C \times B \times C \to D$ defined by
\[
    (b, c, b', c') \mapsto bb' \otimes cc'.
\]
This is $A$-linear in each factor and therefore, by \eqref{prop:multilinear-tensor}, induces an $A$-module homomorphism
\[
    B \otimes C \otimes B \otimes C \to D,
\]
hence by \eqref{pro:canonical-isomorphisms} an $A$-module homomorphism
\[
    D \otimes D \to D
\]
and this in turn by \eqref{prop:additive-function} corresponds to an $A$-bilinear mapping
\[
\mu : D \times D \to D
\]
which is such that
\[
\mu(b \otimes c, b' \otimes c') = bb' \otimes cc'.
\]

Of course, we could have written down this formula directly, but without some such argument as we have given there would be no guarantee that $\mu$ was well-defined.

We have therefore defined a multiplication on the tensor product $D = B \otimes_A C$: for elements of the form $b \otimes c$ it is given by
\[
(b \otimes c)(b' \otimes c') = bb' \otimes cc',
\]
and in general by
\[
\left(\sum_i (b_i \otimes c_i)\right)\left(\sum_j (b'_j \otimes c'_j)\right) = \sum_{i,j} (b_i b'_j \otimes c_i c'_j).
\]

The reader should check that with this multiplication $D$ is a commutative ring, with identity element $1 \otimes 1$. Furthermore, $D$ is an $A$-algebra: the mapping $a \mapsto f(a) \otimes g(a)$ is a ring homomorphism $A \to D$.

In fact there is a commutative diagram of ring homomorphisms
\[
\begin{tikzcd}
 & B \arrow[rd, "u"] & \\
A \arrow[ru, "f"] \arrow[rd, "g"'] & & D \\
 & C \arrow[ru, "v"'] &
\end{tikzcd}
\]
in which $u$, for example, is defined by $u(b) = b \otimes 1$.

%=============================================================================
% Chapter 2 Exercises - With hints from original book
%=============================================================================
\section*{Exercises}
\addcontentsline{toc}{section}{Exercises}

\begin{enumerate}
    \item Show that $(\Z/m\Z) \otimes_\Z (\Z/n\Z) = 0$ if $m, n$ are coprime.

    \item \label{exer2-chapter2}Let $A$ be a ring, $\mathfrak{a}$ an ideal, $M$ an $A$-module. Show that $(A/\mathfrak{a}) \otimes_A M$ is isomorphic to $M/\mathfrak{a}M$.\\
    \textit{Hint:} Tensor the exact sequence $0 \to \mathfrak{a} \to A \to A/\mathfrak{a} \to 0$ with $M$.

    \item \label{exer3-chapter2}Let $A$ be a local ring, $M$ and $N$ finitely generated $A$-modules. Prove that if $M \otimes N = 0$, then $M = 0$ or $N = 0$.\\
    \textit{Hint:} Let $\mathfrak{m}$ be the maximal ideal, $k = A/\mathfrak{m}$ the residue field. Let $M_k = k \otimes_A M \cong M/\mathfrak{m}M$ by Exercise \ref{exer2-chapter2}. By Nakayama's lemma, $M_k = 0 \implies M = 0$. But $M \otimes_A N = 0 \implies (M \otimes_A N)_k = 0 \implies M_k \otimes_k N_k = 0 \implies M_k = 0$ or $N_k = 0$, since $M_k, N_k$ are vector spaces over a field.

    \item \label{exer4-chapter2}Let $M_i$ ($i \in I$) be any family of $A$-modules, and let $M$ be their direct sum. Prove that $M$ is flat $\iff$ each $M_i$ is flat.

    \item \label{exer5-chapter2}Let $A[x]$ be the ring of polynomials in one indeterminate over a ring $A$. Prove that $A[x]$ is a flat $A$-algebra.\\
    \textit{Hint:} Use Exercise \ref{exer4-chapter2}.

    \item \label{exer6-chapter2}For any $A$-module, let $M[x]$ denote the set of all polynomials in $x$ with coefficients in $M$, that is to say expressions of the form $\sum_{i=0}^n m_i x^i$ ($m_i \in M$). Defining the product of an element of $A[x]$ and an element of $M[x]$ in the obvious way, show that $M[x]$ is an $A[x]$-module. 
    
    \quad\, Show that $M[x] \cong A[x] \otimes_A M$.

    \item \label{exer7-chapter2}Let $\mathfrak{p}$ be a prime ideal in $A$. Show that $\mathfrak{p}[x]$ is a prime ideal in $A[x]$. If $\mathfrak{m}$ is a maximal ideal in $A$, is $\mathfrak{m}[x]$ a maximal ideal in $A[x]$?

    \item \label{exer8-chapter2}\begin{enumerate}
        \item If $M$ and $N$ are flat $A$-modules, then so is $M \otimes_A N$.
        \item If $B$ is a flat $A$-algebra and $N$ is a flat $B$-module, then $N$ is flat as an $A$-module.
    \end{enumerate}

    \item \label{exer9-chapter2}Let $0 \to M' \to M \to M'' \to 0$ be an exact sequence of $A$-modules. If $M'$ and $M''$ are finitely generated, then so is $M$.

    \item \label{exer10-chapter2}Let $A$ be a ring $\neq 0$, $\mathfrak{a}$ an ideal contained in the Jacobson radical of $A$; let $M$ be an $A$-module and $N$ a finitely generated $A$-module, and let $u \colon M \to N$ be a homomorphism. If the induced homomorphism $M/\mathfrak{a}M \to N/\mathfrak{a}N$ is surjective, then $u$ is surjective.

    \item \label{exer11-chapter2}Let $A$ be a ring $\neq 0$. Show that $A^m \cong A^n \implies m = n$.\\
    \textit{Hint:} Let $\mathfrak{m}$ be a maximal ideal of $A$ and let $\phi \colon A^m \to A^n$ be an isomorphism. Then $1 \otimes \phi \colon (A/\mathfrak{m}) \otimes A^m \to (A/\mathfrak{m}) \otimes A^n$ is an isomorphism between vector spaces of dimensions $m$ and $n$ over the field $k = A/\mathfrak{m}$. Hence $m = n$. (Cf. Chapter 3, Exercise \ref{exer15-chapter3}.)\\
    \quad\, If $\phi \colon A^m \to A^n$ is surjective, then $m \geq n$.\\
    \quad\, If $\phi \colon A^m \to A^n$ is injective, is it always the case that $m \leq n$?

    \item \label{exer12-chapter2}Let $M$ be a finitely generated $A$-module and $\phi \colon M \to A^n$ a surjective homomorphism. Show that $\Ker(\phi)$ is finitely generated.\\
    \textit{Hint:} Let $e_1, \ldots, e_n$ be a basis of $A^n$ and choose $u_i \in M$ such that $\phi(u_i) = e_i$ ($1 \leq i \leq n$). Show that $M$ is the direct sum of $\Ker(\phi)$ and the submodule generated by $u_1, \ldots, u_n$.

    \item \label{exer13-chapter2}Let $f \colon A \to B$ be a ring homomorphism, and let $N$ be a $B$-module. Regarding $N$ as an $A$-module by restriction of scalars, form the $B$-module $N_B = B \otimes_A N$. Show that the homomorphism $g \colon N \to N_B$ which maps $y$ to $1 \otimes y$ is injective and that $g(N)$ is a direct summand of $N_B$.\\
    \textit{Hint:} Define $p \colon N_B \to N$ by $p(b \otimes y) = by$, and show that $N_B = \Image(g) \oplus \Ker(p)$.

    \emph{Direct limits.}

    \item \label{exer14-chapter2}A partially ordered set $I$ is said to be a \emph{directed set}\index{directed set} if for each pair $i, j$ in $I$ there exists $k \in I$ such that $i \leq k$ and $j \leq k$.\\
    Let $A$ be a ring, let $I$ be a directed set and let $(M_i)_{i \in I}$ be a family of $A$-modules indexed by $I$. For each pair $i, j$ in $I$ such that $i \leq j$, let $\mu_{ij} \colon M_i \to M_j$ be an $A$-homomorphism, and suppose that the following axioms are satisfied:
    \begin{enumerate}[(1)]
        \item $\mu_{ii}$ is the identity mapping of $M_i$, for all $i \in I$;
        \item $\mu_{ik} = \mu_{jk} \circ \mu_{ij}$ whenever $i \leq j \leq k$.
    \end{enumerate}
    Then the modules $M_i$ and homomorphisms $\mu_{ij}$ are said to form a \emph{direct system}\index{direct system} $\mathbf{M} = (M_i, \mu_{ij})$ over the directed set $I$.

    \quad\, We shall construct an $A$-module $M$ called the \emph{direct limit}\index{direct limit} of the direct system $\mathbf{M}$. Let $C$ be the direct sum of the $M_i$, and identify each module $M_i$ with its canonical image in $C$. Let $D$ be the submodule of $C$ generated by all elements of the form $x_i - \mu_{ij}(x_i)$ where $i \leq j$ and $x_i \in M_i$. Let $M = C/D$, let $\mu \colon C \to M$ be the projection and let $\mu_i$ be the restriction of $\mu$ to $M_i$.

    \quad\, The module $M$, or more correctly the pair consisting of $M$ and the family of homomorphisms $\mu_i \colon M_i \to M$, is called the direct limit of the direct system $\mathbf{M}$, and is written $\varinjlim M_i$. From the construction it is clear that $\mu_i = \mu_j \circ \mu_{ij}$ whenever $i \leq j$.

    \item \label{exer15-chapter2}In the situation of Exercise 14, show that every element of $M$ can be written in the form $\mu_i(x_i)$ for some $i \in I$ and some $x_i \in M_i$. Show also that if $\mu_i(x_i) = 0$ then there exists $j \geq i$ such that $\mu_{ij}(x_i) = 0$ in $M_j$.

    \item \label{exer16-chapter2}Show that the direct limit is characterized (up to isomorphism) by the following property. Let $N$ be an $A$-module and for each $i \in I$ let $\alpha_i \colon M_i \to N$ be an $A$-module homomorphism such that $\alpha_i = \alpha_j \circ \mu_{ij}$ whenever $i \leq j$. Then there exists a unique homomorphism $\alpha \colon M \to N$ such that $\alpha_i = \alpha \circ \mu_i$ for all $i \in I$.

    \item \label{exer17-chapter2}Let $(M_i)_{i \in I}$ be a family of submodules of an $A$-module, such that for each pair of indices $i, j$ in $I$ there exists $k \in I$ such that $M_i + M_j \subseteq M_k$. Define $i \leq j$ to mean $M_i \subseteq M_j$ and let $\mu_{ij} \colon M_i \to M_j$ be the embedding of $M_i$ in $M_j$. Show that
    \[
    \varinjlim M_i = \sum M_i = \bigcup M_i.
    \]
    In particular, any $A$-module is the direct limit of its finitely generated submodules.

    \item \label{exer18-chapter2}Let $\mathbf{M} = (M_i, \mu_{ij})$, $\mathbf{N} = (N_i, \nu_{ij})$ be direct systems of $A$-modules over the same directed set. Let $M, N$ be the direct limits and $\mu_i \colon M_i \to M$, $\nu_i \colon N_i \to N$ the associated homomorphisms. 
    
    A \emph{homomorphism} $\Phi \colon \mathbf{M} \to \mathbf{N}$ is by definition a family of $A$-module homomorphisms $\phi_i \colon M_i \to N_i$ such that $\phi_j \circ \mu_{ij} = \nu_{ij} \circ \phi_i$ whenever $i \leq j$. Show that $\Phi$ defines a unique homomorphism $\phi = \varinjlim \phi_i \colon M \to N$ such that $\phi \circ \mu_i = \nu_i \circ \phi_i$ for all $i \in I$.

    \item \label{exer19-chapter2}A sequence of direct systems and homomorphisms $\mathbf{M} \to \mathbf{N} \to \mathbf{P}$ is exact if the corresponding sequence of modules and module homomorphisms is exact for each $i \in I$. Show that the sequence $\varinjlim \mathbf{M} \to \varinjlim \mathbf{N} \to \varinjlim \mathbf{P}$ of direct limits is then exact.\\
    \textit{Hint:} Use Exercise \ref{exer15-chapter2}.

    \emph{Tensor products commute with direct limits.} 

    \item \label{exer20-chapter2}Keeping the same notation as in Exercise 14, let $N$ be any $A$-module. Then $(M_i \otimes N, \mu_{ij} \otimes 1)$ is a direct system; let $P = \varinjlim (M_i \otimes N)$ be its direct limit.
    
    For each $i \in I$ we have a homomorphism $\mu_i \otimes 1 \colon M_i \otimes N \to M \otimes N$, hence by Exercise 16 a homomorphism $\psi \colon P \to M \otimes N$. Show that $\psi$ is an isomorphism, so that
    \[
    \varinjlim (M_i \otimes N) \cong (\varinjlim M_i) \otimes N.
    \]
    \textit{Hint:} For each $i \in I$, let $g_i \colon M_i \times N \to M_i \otimes N$ be the canonical bilinear mapping. Passing to the limit we obtain a mapping $g \colon M \times N \to P$. Show that $g$ is $A$-bilinear and hence define a homomorphism $\phi \colon M \otimes N \to P$. Verify that $\phi \circ \psi$ and $\psi \circ \phi$ are identity mappings.

    \item \label{exer21-chapter2}Let $(A_i)_{i \in I}$ be a family of rings indexed by a directed set $I$, and for each pair $i \leq j$ in $I$ let $\alpha_{ij} \colon A_i \to A_j$ be a ring homomorphism, satisfying conditions (1) and (2) of Exercise \ref{exer14-chapter2}. Regarding each $A_i$ as a $\Z$-module we can then form the direct limit $A = \varinjlim A_i$. Show that $A$ inherits a ring structure from the $A_i$ so that the mappings $A_i \to A$ are ring homomorphisms. The ring $A$ is the \emph{direct limit} of the system $(A_i, \alpha_{ij})$.
    
    \quad\, If $A = 0$ prove that $A_i = 0$ for some $i \in I$.

    \textit{Hint:} Remember that all rings have identity elements!

    \item \label{exer22-chapter2}Let $(A_i, \alpha_{ij})$ be a direct system of rings and let $\mathfrak{N}_i$ be the nilradical of $A_i$. Show that $\varinjlim \mathfrak{N}_i$ is the nilradical of $\varinjlim A_i$.
    
    If each $A_i$ is an integral domain, then $\varinjlim A_i$ is an integral domain.

    \item \label{exer23-chapter2}Let $(B_\lambda)_{\lambda \in \Lambda}$ be a family of $A$-algebras. For each finite subset of $\Lambda$ let $B_J$ denote the tensor product (over $A$) of the $B_\lambda$ for $\lambda \in J$. If $J'$ is another finite subset of $\Lambda$ and $J \subseteq J'$, there is a canonical $A$-algebra homomorphism $B_J \to B_{J'}$. Let $B$ denote the direct limit of the rings $B_J$ as $J$ runs through all finite subsets of $\Lambda$. The ring $B$ has a natural $A$-algebra structure for which the homomorphisms $B_J \to B$ are $A$-algebra homomorphisms. The $A$-algebra $B$ is the \emph{tensor product}\index{tensor product} of the family $(B_\lambda)_{\lambda \in \Lambda}$.

    \emph{Flatness and Tor.} 

    In these Exercises it will be assumed that the reader is familiar with the definition and basic properties of the Tor functor.

    \item \label{exer24-chapter2}If $M$ is an $A$-module, the following are equivalent:
    \begin{enumerate}[i)]
        \item $M$ is flat;
        \item $\Tor_n^A(M, N) = 0$ for all $n > 0$ and all $A$-modules $N$;
        \item $\Tor_1^A(M, N) = 0$ for all $A$-modules $N$.
    \end{enumerate}
    \textit{Hint:} To show that (i) $\implies$ (ii), take a free resolution of $N$ and tensor it with $M$. Since $M$ is flat, the resulting sequence is exact and therefore its homology groups, which are the $\Tor_n^A(M, N)$, are zero for $n > 0$. To show that (iii) $\implies$ (i), let $0 \to N' \to N \to N'' \to 0$ be an exact sequence. Then, from the Tor exact sequence,
    \[
    \Tor_1^A(M, N'') \to M \otimes N' \to M \otimes N \to M \otimes N'' \to 0
    \]
    is exact. Since $\Tor_1^A(M, N'') = 0$ it follows that $M$ is flat.

    \item \label{exer25-chapter2}Let $0 \to N' \to N \to N'' \to 0$ be an exact sequence, with $N''$ flat. Then $N'$ is flat $\iff$ $N$ is flat.

    \textit{Hint:} Use Exercise \ref{exer24-chapter2} and the Tor exact sequence.

    \item \label{exer26-chapter2}Let $N$ be an $A$-module. Then $N$ is flat $\iff$ $\Tor_1^A(A/\mathfrak{a}, N) = 0$ for all finitely generated ideals $\mathfrak{a}$ in $A$.\\
    \textit{Hint:} Show first that $N$ is flat if $\Tor_1^A(M, N) = 0$ for all finitely generated $A$-modules $M$, by using \eqref{prop:flat-equivalent}. If $M$ is finitely generated, let $x_1, \ldots, x_n$ be a set of generators of $M$, and let $M_i$ be the submodule generated by $x_1, \ldots, x_i$. By considering the successive quotients $M_i/M_{i-1}$ and using Exercise \ref{exer25-chapter2}, deduce that $N$ is flat if $\Tor_1^A(M, N) = 0$ for all \emph{cyclic}\index{cyclic} $A$-modules $M$, i.e., all $M$ generated by a single element, and therefore of the form $A/\mathfrak{a}$ for some ideal $\mathfrak{a}$. Finally use \eqref{prop:flat-equivalent} again to reduce to the case where $\mathfrak{a}$ is a finitely generated ideal.

    \item \label{exer27-chapter2}A ring $A$ is \emph{absolutely flat}\index{absolutely flat ring} if every $A$-module is flat. Prove that the following are equivalent:
    \begin{enumerate}[i)]
        \item $A$ is absolutely flat.
        \item Every principal ideal is idempotent.
        \item Every finitely generated ideal is a direct summand of $A$.
    \end{enumerate}
    \textit{Hint:} (i) $\implies$ (ii). Let $x \in A$. Then $A/(x)$ is a flat $A$-module, hence in the diagram
    \[\begin{tikzcd}
        {(x)\otimes A} & {(x)\otimes A/(x)} \\
        A & {A/(x)}
        \arrow["\beta", from=1-1, to=1-2]
        \arrow[from=1-1, to=2-1]
        \arrow["\alpha", from=1-2, to=2-2]
        \arrow[from=2-1, to=2-2]
    \end{tikzcd}\]
    the mapping $\alpha$ is injective. Hence $\Image(\beta) = 0$, hence $(x) = (x^2)$.\\
    (ii) $\implies$ (iii). Let $x \in A$. Then $x = ax^2$ for some $a \in A$, hence $e = ax$ is idempotent and we have $(e) = (x)$. Now if $e, f$ are idempotents, then $(e, f) = (e + f - ef)$. Hence every finitely generated ideal is principal, and generated by an idempotent $e$, hence is a direct summand because $A = (e) \oplus (1-e)$.\\
    (iii) $\implies$ (i). Use the criterion of Exercise \ref{exer26-chapter2}.

    \item \label{exer28-chapter2}A Boolean ring is absolutely flat. The ring of Chapter 1, Exercise \ref{exer7-chapter1} is absolutely flat. Every homomorphic image of an absolutely flat ring is absolutely flat. If a local ring is absolutely flat, then it is a field.
    
    \quad\, If $A$ is absolutely flat, every non-unit in $A$ is a zero-divisor.
\end{enumerate}

